%======================================================================
%  PulseIQ — Technical Documentation
%  Auto‑generated LaTeX source
%======================================================================
\documentclass[12pt, a4paper, oneside]{report}

%----------------------------------------------------------------------
% Packages
%----------------------------------------------------------------------
\usepackage[utf8]{inputenc}
\usepackage[T1]{fontenc}
\usepackage{lmodern}
\usepackage[margin=1in]{geometry}
\usepackage{graphicx}
\usepackage{xcolor}
\usepackage{hyperref}
\usepackage{listings}
\usepackage{booktabs}
\usepackage{enumitem}
\usepackage{amsmath, amssymb}
\usepackage{fancyhdr}
\usepackage{titlesec}
\usepackage{tocloft}
\usepackage{parskip}
\usepackage{microtype}
\usepackage{caption}
\usepackage{float}

%----------------------------------------------------------------------
% Colour palette
%----------------------------------------------------------------------
\definecolor{pulseblue}{HTML}{1A73E8}
\definecolor{pulsedark}{HTML}{1B1F3B}
\definecolor{pulsegray}{HTML}{5F6368}
\definecolor{codebackground}{HTML}{F5F5F5}
\definecolor{codekeyword}{HTML}{0D47A1}
\definecolor{codestring}{HTML}{1B5E20}
\definecolor{codecomment}{HTML}{9E9E9E}

%----------------------------------------------------------------------
% Hyperref configuration
%----------------------------------------------------------------------
\hypersetup{
    colorlinks  = true,
    linkcolor   = pulseblue,
    urlcolor    = pulseblue,
    citecolor   = pulseblue,
    pdfauthor   = {PulseIQ Team},
    pdftitle    = {PulseIQ Technical Documentation},
    pdfsubject  = {Employee Burnout Detection Platform},
}

%----------------------------------------------------------------------
% Listings style
%----------------------------------------------------------------------
\lstdefinestyle{pulseCode}{
    backgroundcolor  = \color{codebackground},
    basicstyle       = \ttfamily\small,
    keywordstyle     = \color{codekeyword}\bfseries,
    stringstyle      = \color{codestring},
    commentstyle     = \color{codecomment}\itshape,
    breaklines       = true,
    frame            = single,
    rulecolor        = \color{pulsegray!40},
    numbers          = left,
    numberstyle      = \tiny\color{pulsegray},
    numbersep        = 8pt,
    xleftmargin      = 1.5em,
    framexleftmargin = 1.5em,
    tabsize          = 4,
    showstringspaces = false,
    captionpos       = b,
}
\lstset{style=pulseCode}

%----------------------------------------------------------------------
% Header / Footer
%----------------------------------------------------------------------
\pagestyle{fancy}
\fancyhf{}
\fancyhead[L]{\small\textcolor{pulsegray}{\leftmark}}
\fancyhead[R]{\small\textcolor{pulsegray}{PulseIQ Documentation}}
\fancyfoot[C]{\thepage}
\renewcommand{\headrulewidth}{0.4pt}
\renewcommand{\footrulewidth}{0pt}

%----------------------------------------------------------------------
% Chapter / section formatting
%----------------------------------------------------------------------
\titleformat{\chapter}[display]
  {\normalfont\Huge\bfseries\color{pulsedark}}
  {\chaptertitlename\ \thechapter}{20pt}{\Huge}
\titleformat{\section}
  {\normalfont\Large\bfseries\color{pulsedark}}{\thesection}{1em}{}
\titleformat{\subsection}
  {\normalfont\large\bfseries\color{pulsedark}}{\thesubsection}{1em}{}
\titleformat{\subsubsection}
  {\normalfont\normalsize\bfseries\color{pulsedark}}{\thesubsubsection}{1em}{}

%----------------------------------------------------------------------
% Table of Contents styling
%----------------------------------------------------------------------
\renewcommand{\cftchapfont}{\bfseries\color{pulsedark}}
\renewcommand{\cftsecfont}{\color{pulsedark}}

%======================================================================
%  DOCUMENT START
%======================================================================
\begin{document}

%----------------------------------------------------------------------
%  Title Page
%----------------------------------------------------------------------
\begin{titlepage}
    \centering
    \vspace*{2cm}

    {\Huge\bfseries\textcolor{pulsedark}{PulseIQ}}\\[0.5cm]
    {\LARGE\textcolor{pulseblue}{Technical Documentation}}\\[2cm]

    {\large\textcolor{pulsegray}{%
        Advanced Employee Burnout Detection, Tracking,\\
        and Mitigation Platform%
    }}\\[3cm]

    \rule{0.6\textwidth}{0.5pt}\\[1cm]

    {\large
        \textbf{Architecture} $\bullet$
        \textbf{Pipeline} $\bullet$
        \textbf{Algorithms} $\bullet$
        \textbf{Machine Learning} $\bullet$
        \textbf{Gen AI} $\bullet$
        \textbf{Integration}
    }\\[2cm]

    {\large\textcolor{pulsegray}{Version 2.0 --- Now with Generative AI}}\\[0.5cm]
    {\large\textcolor{pulsegray}{\today}}

    \vfill
\end{titlepage}

%----------------------------------------------------------------------
%  Table of Contents
%----------------------------------------------------------------------
\tableofcontents
\newpage

%======================================================================
%  CHAPTER 1 — Overview & Running Instructions
%======================================================================
\chapter{Project Overview \& Running Instructions}
\label{ch:overview}

PulseIQ is an advanced full-stack application designed to \textbf{detect}, \textbf{track}, and \textbf{mitigate} employee burnout using a suite of statistical, behavioral, and machine learning models.

%----------------------------------------------------------------------
\section{High-Level Architecture}
\label{sec:architecture}

The system comprises three major tiers:

\begin{enumerate}[label=\textbf{\arabic*.}]
    \item \textbf{Machine Learning Pipeline (Python)}\\
        Generates or processes synthetic workspace data across text, meetings, and code repositories to derive a cohesive, ensemble-driven \emph{Burnout Probability} per employee.

    \item \textbf{Backend API (FastAPI)}\\
        Serves the resulting metrics, time-series projections, ML predictions, and AI-driven narratives to the client.

    \item \textbf{Frontend Dashboard (React / Vite)}\\
        Gives employees and managers detailed insights, alerts, and actionable recommendations.
\end{enumerate}

%----------------------------------------------------------------------
\section{How to Run the Program}
\label{sec:running}

The application is bundled into a unified launcher \texttt{run.py} at the project root which coordinates the ML processing and both web environments simultaneously.

\subsection{Prerequisites}

Before running, ensure the following are installed:

\begin{itemize}
    \item Python 3.9+
    \item Node.js v16+
    \item All Python and JavaScript dependencies:
\end{itemize}

\begin{lstlisting}[language=bash, caption={Installing dependencies}]
pip install -r backend/requirements.txt
npm install
\end{lstlisting}

\subsection{General Startup}

To run the full suite (regenerate all ML data, start API, start React):

\begin{lstlisting}[language=bash, caption={Full application launch}]
python run.py
\end{lstlisting}

\subsection{Advanced Modes}

\texttt{run.py} supports several configurations depending on your needs:

\begin{table}[H]
\centering
\caption{Command-line flags for \texttt{run.py}}
\label{tab:flags}
\begin{tabular}{@{} l p{8.5cm} @{}}
\toprule
\textbf{Flag} & \textbf{Description} \\
\midrule
\texttt{-{}-skip-pipeline} & Skips the ML pipeline; starts only the FastAPI and Vite servers. Use when \texttt{pulseiq\_data/} is already up to date.\\[4pt]
\texttt{-{}-skip-ml}       & Runs the standard behavioral pipeline \emph{without} the advanced ML models (LSTM, time-series, anomalies). Faster execution.\\[4pt]
\texttt{-{}-api-only}      & Starts only the FastAPI server (\texttt{http://localhost:8000}).\\[4pt]
\texttt{-{}-frontend-only} & Starts only the Vite UI server (\texttt{http://localhost:5173}).\\
\bottomrule
\end{tabular}
\end{table}

\subsection{What Happens During \texttt{python run.py}?}

\begin{enumerate}
    \item The script first runs \texttt{backend/main.py}. This orchestrates 10 sequential pipeline steps:
    \begin{center}
        Data Generation $\rightarrow$ Feature Building $\rightarrow$ Model Training $\rightarrow$ Scoring $\rightarrow$ Ensemble Generation $\rightarrow$ Final Reporting
    \end{center}
    Data artifacts are generated in \texttt{pulseiq\_data/}.

    \item Upon successful pipeline completion, the system daemonizes an API server (\texttt{uvicorn}) and a development frontend server (\texttt{npm run dev}).

    \item Pressing \texttt{Ctrl+C} cleanly tears down the application using automated PID trapping.
\end{enumerate}


%======================================================================
%  CHAPTER 2 — Pipeline Orchestration
%======================================================================
\chapter{Pipeline Orchestration}
\label{ch:pipeline}

The ML pipeline is contained within \texttt{backend/main.py}, which sequences together discrete logic chunks into a deterministic execution flow.

%----------------------------------------------------------------------
\section{The STEPS Framework}
\label{sec:steps}

\texttt{main.py} registers an execution framework structured as an array of dictionaries (the \texttt{STEPS} array). It dynamically imports and invokes each module's \texttt{main()} function, gracefully catching errors at every stage.

The deterministic execution sequence is:

\begin{enumerate}[label=\textbf{Step \arabic*.}]
    \item \texttt{generate\_data.py} --- Constructs the synthetic ground-truth ecosystem.
    \item \texttt{aggregate\_features.py} --- Groups metrics by \texttt{(employee, date)} pairs for ML compatibility.
    \item \texttt{compute\_baselines.py} --- Derives rolling personal benchmarks (what does ``normal'' mean for \emph{this} employee).
    \item \texttt{compute\_burnout.py} --- The behavioral mathematical engine evaluating rest decay and cognitive overload.
    \item \texttt{recommendations.py} --- Assigns tailored health guidance and alerts managers.
    \item \texttt{manager\_insights.py} --- Pre-aggregates rollups for team-scale charting.
    \item \texttt{run\_ml\_parallel.py} --- Kicks off the parallel deep learning jobs (Time-Series ARIMA, Isolation Forests for anomalous drops, PyTorch/TensorFlow LSTMs).
    \item \texttt{predictive\_ensemble.py} --- Weighs all inputs to form an ultimate confidence-rated diagnosis.
    \item \texttt{generate\_report.py} --- Builds a cohesive terminal summary of the riskiest individuals.
\end{enumerate}

%----------------------------------------------------------------------
\section{Logging and Fault Tolerance}
\label{sec:logging}

Each step logs its status, execution elapsed time, and output object scale. The orchestrator employs a two-tier fault tolerance strategy:

\begin{itemize}
    \item \textbf{Core process failure:} If a fundamental pipeline step fails (e.g., data generation, feature engineering), \texttt{main.py} gracefully terminates the entire pipeline.

    \item \textbf{Advanced ML subprocess failure:} If an advanced ML process fails (e.g., an LSTM batch taking excessive VRAM), the orchestrator yields a warning but permits the pipeline to continue onto the Predictive Ensemble with available data.
\end{itemize}


%======================================================================
%  CHAPTER 3 — Data Generation & Feature Engineering
%======================================================================
\chapter{Data Generation \& Feature Engineering}
\label{ch:data}

Because PulseIQ models complex workplace interactions, raw ingestion mimics realistic signals an HR portal would extract from SaaS suites.

%----------------------------------------------------------------------
\section{Synthetic Data Generation}
\label{sec:synth}

\texttt{generate\_data.py} generates pseudo-random events over several weeks representing realistic workflow signals:

\begin{table}[H]
\centering
\caption{Synthetic data sources and their dimensions}
\label{tab:datasources}
\begin{tabular}{@{} l p{9cm} @{}}
\toprule
\textbf{Source} & \textbf{What is Modeled} \\
\midrule
\textbf{Slack}     & Raw messaging habits---mapped as Direct Messages vs.\ Work Channels vs.\ Social Channels. \\[4pt]
\textbf{Jira}      & Ticket comments and velocity by calculating \texttt{story\_points} assigned $\rightarrow$ resolved. \\[4pt]
\textbf{Calendar}  & Day structures including start times and attendee counts. \\[4pt]
\textbf{Git}       & Line additions, deletions, repo interactions, with high likelihood of weekend or after-hours commits for ``stressed'' patterns. \\[4pt]
\textbf{Zoom}      & Call duration and active speaker engagement. \\
\bottomrule
\end{tabular}
\end{table}

\textbf{Output:} \texttt{pulseiq\_data/*\_events.csv} (and similar artifacts).

%----------------------------------------------------------------------
\section{Feature Engineering}
\label{sec:featureeng}

\texttt{aggregate\_features.py} translates the raw chronological payloads into a structured standard where each record corresponds to a unique \texttt{[Employee, Date]} pair.

\subsection{After-Hours Filtering}

Employs an explicit temporal boundary (\textbf{21:00--07:00}) for Git and Slack data to attribute the \texttt{after\_hours\_ratio} feature, which captures the proportion of activity occurring outside normal working hours.

\subsection{RoBERTa NLP Integration}

Leverages the Hugging Face \texttt{cardiffnlp/twitter-roberta-base-sentiment} pipeline:

\begin{itemize}
    \item Messages are interpreted securely, turning text into floats:
    \[
        \text{sentiment} \in [0.0, 1.0], \quad 1.0 = \text{highly positive}, \quad 0.0 = \text{highly negative}
    \]

    \item \textbf{Fallback mechanism:} If the local machine lacks PyTorch or Transformers, the system leverages existing synthetic baseline scores.

    \item Extrapolates NLP arrays over a day into:
    \begin{itemize}
        \item \texttt{slack\_sentiment\_volatility} --- mathematical standard deviation ($\sigma$) of mood variance.
        \item \texttt{slack\_avg\_stress\_score} --- fractional occurrence of trigger words such as \emph{``overloaded''}, \emph{``swamped''}, \emph{``drowning''}.
    \end{itemize}
\end{itemize}

\subsection{Meeting Concentration (Largest Focus Gap)}

Sorts meeting interval arrays per employee, then iterates and calculates the longest continuous block spanning their schedule to measure how fragmented a workday truly is.

\subsection{Zoom Engagement Ratios}

Computes \texttt{zoom\_speaking\_ratio}:
\[
    \texttt{zoom\_speaking\_ratio} = \frac{\sum \texttt{speaking\_seconds}}{\sum \texttt{duration}}
\]

This ratio acts as a potent model proxy for social withdrawal or disconnection from collaborative activities.

\textbf{Output:} \texttt{pulseiq\_data/daily\_features.csv}


%======================================================================
%  CHAPTER 4 — Behavioral Burnout Engine
%======================================================================
\chapter{Behavioral Burnout Engine}
\label{ch:burnout}

\texttt{compute\_burnout.py} is the bedrock behavioral analysis system. It maps statistical deviations of employee workflow against their personal baselines to model \emph{cognitive load}, \emph{rest deficit}, and a \emph{raw burnout scale}.

%----------------------------------------------------------------------
\section{Normalization Against Baselines}
\label{sec:baselines}

Each employee's history has a unique baseline computed by \texttt{compute\_baselines.py} using robust statistics. For instance, 10~meetings per day may be standard for product managers but highly anomalous for engineers.

When calculating risk, metrics are normalized using a \textbf{Robust Z-Score} against the 14-day trailing behavior:

\begin{lstlisting}[language=Python, caption={Robust Z-Score normalization}]
def robust_z(value, median, std_robust):
    """
    Robust Z-Score logic avoiding zero-division
    on uniform behavior.
    """
    if std_robust == 0:
        return 0.0
    return (value - median) / std_robust
\end{lstlisting}

Mathematically:
\[
    z_{\text{robust}} = \frac{x - \tilde{x}}{\hat{\sigma}_{\text{robust}}}
\]

where $\tilde{x}$ is the trailing median and $\hat{\sigma}_{\text{robust}}$ is the robust standard deviation over the past 14~days.

%----------------------------------------------------------------------
\section{The Three Pillars of Burnout}
\label{sec:pillars}

The behavioral engine decomposes burnout into three orthogonal sub-scores, each ranging from 0 to 100:

\subsection{Deep Work Index (0--100)}

Assesses meaningful output generation by evaluating:
\begin{itemize}
    \item Jira Story Points closed
    \item Git code diffs (net lines changed)
    \item Largest Focus Gap Minutes from calendar analysis
\end{itemize}

Each factor is evaluated as a ratio relative to the employee's benchmark metric.

\subsection{Fragmentation Score (0--100)}

High fragmentation destroys cognitive momentum. This score aggregates:
\begin{itemize}
    \item Meeting count
    \item Slack total messages
    \item Zoom meeting frequency
\end{itemize}

Z-scores are mapped through a non-linear $\tanh$ function, pushing maximum distraction asymptotically towards 100:
\[
    \text{Fragmentation} = 100 \cdot \tanh\!\left(\alpha \cdot z_{\text{agg}}\right)
\]

\subsection{Connection Index (0--100)}

Measures workplace social isolation. Driven primarily by:
\begin{itemize}
    \item NLP-derived Slack sentiment ratios
    \item Distinct DM partner count (peer engagement vs.\ isolation)
    \item Zoom speaking ratios
\end{itemize}

%----------------------------------------------------------------------
\section{Rest Deficit \& Trailing Probabilities}
\label{sec:restdeficit}

\begin{enumerate}
    \item \textbf{Consecutive Work Days:} Tracks unrelenting strings of active contribution (after-hours commits, Jira actions on weekends).

    \item \textbf{Composite Score:} Combines the three pillars against the Rest Deficit and Recovery histories to form a raw mathematical score.

    \item \textbf{Sigmoid Transformation:} Converts the score to a probability in $[0, 1]$:
    \[
        P(\text{burnout}) = \sigma(s) = \frac{1}{1 + e^{-s}}
    \]

    \item \textbf{3-Day Simple Moving Average (SMA):} Smooths the probability so a single bad day does not misclassify someone as burnt out:
    \[
        \text{SMA}_3(t) = \frac{P(t) + P(t-1) + P(t-2)}{3}
    \]

    \item \textbf{Forecast to Critical:} Generates a linear slope-based projection line \texttt{forecast\_days\_until\_critical}, extrapolating when the employee's SMA will hit the $0.85$ critical zone.
\end{enumerate}


%======================================================================
%  CHAPTER 5 — Machine Learning Suite
%======================================================================
\chapter{Machine Learning Suite}
\label{ch:ml}

The PulseIQ backend supplements the core behavioral engine with advanced time-series analysis and neural networks. Execution is coordinated by \texttt{run\_ml\_parallel.py}, which delegates individual modeling routines.

%----------------------------------------------------------------------
\section{Time-Series Analysis}
\label{sec:timeseries}

\texttt{time\_series\_analysis.py} treats the trailing probabilities of burnout as a sequential mathematical wave.

\begin{itemize}
    \item \textbf{EWMA (Exponentially Weighted Moving Average):} Highlights the macro-signal over the noise:
    \[
        \text{EWMA}_t = \alpha \cdot x_t + (1 - \alpha) \cdot \text{EWMA}_{t-1}
    \]

    \item \textbf{Forecasting:} Identifies the underlying behavioral distribution curve to generate forward-looking projections.

    \item \textbf{Rolling Volatility:} Signifies erratic schedules or unstable work-life borders.

    \item \textbf{Trend Direction:} Classifies current trajectory as \texttt{rising\_fast}, \texttt{rising}, \texttt{stable}, \texttt{declining}, etc.
\end{itemize}

\textbf{Output:} \texttt{pulseiq\_data/timeseries\_forecasts.csv}

%----------------------------------------------------------------------
\section{Anomaly Detection}
\label{sec:anomaly}

\texttt{anomaly\_detection.py} utilizes Scikit-Learn's \texttt{IsolationForest} to monitor multidimensional anomalies across dozens of variables simultaneously.

\begin{itemize}
    \item While the behavioral engine checks against an \emph{absolute} scale, the Isolation Forest maps employees \emph{dimensionally}.
    \item It tags an event as anomalous if the employee deviates substantially from their localized clusters---for example, if they suddenly drop off Jira, shift code commitments rapidly, and withdraw from DMs overnight.
\end{itemize}

\textbf{Output:} \texttt{pulseiq\_data/anomaly\_scores.csv}, tagging exact outlier behaviors.

%----------------------------------------------------------------------
\section{Deep Learning Prediction (LSTM)}
\label{sec:lstm}

\texttt{deep\_learning\_model.py} deploys \textbf{Long Short-Term Memory} (LSTM) networks built using TensorFlow/PyTorch architectures.

\begin{itemize}
    \item LSTMs inherently retain internal cell state over multiple observations (the ``trailing sequence barrier''), enabling them to capture long-range temporal dependencies.

    \item \textbf{Input:} Sequences of trailing 14--30 day \texttt{daily\_features} + \texttt{daily\_scores} data matrices.

    \item \textbf{Capability:} Recognizes nonlinear patterns that signify impending drops \emph{before} traditional standard-deviation triggers fire.
\end{itemize}

\textbf{Output:} \texttt{pulseiq\_data/lstm\_predictions.csv}, carrying its own confidence mapping.


%======================================================================
%  CHAPTER 6 — Assembly & Integration
%======================================================================
\chapter{Assembly \& Integration Flow}
\label{ch:assembly}

Once the ML modeling logic concludes, PulseIQ must combine disparate findings, generate textual warnings, and format the data accurately for web transport.

%----------------------------------------------------------------------
\section{Predictive Ensemble}
\label{sec:ensemble}

\texttt{predictive\_ensemble.py} combines predictions into a unified truth. Since ML models generate varying levels of hallucination or variance, they require stabilization.

\subsection{Weighted Average Framework}

Models combine along a weighted average confidence framework based on model reliability:

\begin{table}[H]
\centering
\caption{Ensemble model weights}
\label{tab:weights}
\begin{tabular}{@{} l c p{6.5cm} @{}}
\toprule
\textbf{Model} & \textbf{Weight} & \textbf{Rationale} \\
\midrule
Behavioral Engine     & 35\% & Deterministic engine---highest stable accuracy. \\[3pt]
Time Series (ARIMA)   & 20\% & Moderate weight, strong pattern identifier. \\[3pt]
LSTM Neural Network   & $30\% \times C_{\text{dyn}}$ & High weight but actively scales itself downward if the neural network is uncertain ($C_{\text{dyn}} \in [0,1]$). \\[3pt]
Anomaly Detection     & Multiplier & If anomalous behavior triggers, amplifies underlying score estimates. \\
\bottomrule
\end{tabular}
\end{table}

\subsection{Disagreement Index}

The ensemble automatically determines the \textbf{Disagreement Index} when models conflict (e.g., ARIMA predicts healthy but LSTM predicts disaster), which yields the final model confidence interval.

\subsection{Risk Tier Classification}

Probabilities are grouped into five risk tiers:

\begin{center}
\begin{tabular}{@{} l l @{}}
\toprule
\textbf{Tier} & \textbf{Probability Range} \\
\midrule
\texttt{CRITICAL} & $\geq 0.85$ \\
\texttt{HIGH}     & $[0.65, 0.85)$ \\
\texttt{MEDIUM}   & $[0.45, 0.65)$ \\
\texttt{LOW}      & $[0.25, 0.45)$ \\
\texttt{MINIMAL}  & $< 0.25$ \\
\bottomrule
\end{tabular}
\end{center}

\subsection{Narrative Generation}

The ensemble automatically composes a sentence-based narrative using rule-based parsing logic, for example:
\begin{quote}
    \emph{``High fragmentation from excessive meetings. A behavioral change-point was recently detected.''}
\end{quote}

%----------------------------------------------------------------------
\section{Rule-Based Recommendations}
\label{sec:recommendations}

\texttt{recommendations.py} applies deterministic mitigation strategies mapped from resulting data flags:

\begin{itemize}
    \item \textbf{High Interruptions detected} $\rightarrow$ Schedule Focus Blocks
    \item \textbf{After-hours pushes flagged} $\rightarrow$ Enforce Work-Life Boundary
    \item \textbf{Workload imbalance} $\rightarrow$ Workload Reallocation (peer matching heuristic)
\end{itemize}

%----------------------------------------------------------------------
\section{API Server}
\label{sec:api}

The Python \textbf{FastAPI} application utilizing \texttt{Uvicorn} serves as the data bridge:

\begin{itemize}
    \item Runs on \texttt{http://localhost:8000}
    \item Configures generic CORS logic for cross-origin frontend access
    \item Serves data endpoints:
    \begin{itemize}
        \item \texttt{/api/employees} --- Employee listing and metrics
        \item \texttt{/api/health} --- Health check
        \item \texttt{/api/suggestions} --- AI-driven recommendations
        \item \texttt{/api/employees/\{id\}/history} --- Historical time-series data
        \item \texttt{/api/ensemble/summary} --- Ensemble risk distribution
    \end{itemize}
    \item Serves Generative AI (LLM) endpoints powered by Gemini Flash (see Chapter~\ref{ch:genai}):
    \begin{itemize}
        \item \texttt{/api/llm/draft\_message} --- AI-drafted manager messages
        \item \texttt{/api/llm/narrative/\{id\}} --- Dynamic employee narratives
        \item \texttt{/api/llm/ask\_pulse} --- Natural language copilot queries
    \end{itemize}
    \item Loads resultant CSV artifacts directly from \texttt{pulseiq\_data/} and shapes them into normalized JSON dictionaries that the React TypeScript application natively ingests
    \item Formats dates, safely assigns Manager structures, handles \texttt{NaN} anomalies gracefully, and builds historical time-series graphing arrays
\end{itemize}

%----------------------------------------------------------------------
\section{Frontend Design}
\label{sec:frontend}

The React application (\texttt{src/}) binds to the FastAPI interfaces:

\begin{itemize}
    \item Built as a \textbf{Vite} application supporting auto-reload (\texttt{npm run dev})
    \item \textbf{Manager Dashboards:} Maps directly correlated reports for team overview
    \item \textbf{Employee Views:} Localized, personalized burnout insights
    \item Integrates \texttt{recharts} to chart historical tracking statistics supplied by the aggregated \texttt{api.py} payloads
\end{itemize}


%======================================================================
%  CHAPTER 7 — Generative AI Features (Gemini Flash)
%======================================================================
\chapter{Generative AI Features (Gemini Flash)}
\label{ch:genai}

PulseIQ extends its intelligence layer beyond classical ML and statistical models with \textbf{Generative AI} capabilities powered by \textbf{Google Gemini 2.5 Flash}. These features complement the deterministic pipeline by introducing natural language understanding and generation, enabling managers to interact with burnout data conversationally and communicate interventions empathetically.

The Gen AI subsystem is implemented as a set of LLM-backed API endpoints within \texttt{backend/api.py}, utilizing the \texttt{google-genai} Python SDK to interface with the Gemini model.

%----------------------------------------------------------------------
\section{Architecture Overview}
\label{sec:genai-arch}

All Gen AI features follow a consistent pattern:

\begin{enumerate}
    \item \textbf{Context Assembly:} The API gathers relevant employee metrics, team data, and/or recommendation outputs from the internal data layer.
    \item \textbf{Prompt Engineering:} A carefully structured prompt is constructed with role instructions, data context, and task-specific constraints (tone, format, length).
    \item \textbf{LLM Invocation:} The prompt is sent to \texttt{gemini-2.5-flash} via the \texttt{google.genai} client.
    \item \textbf{Response Delivery:} The generated text is returned to the frontend as a JSON payload.
\end{enumerate}

\begin{lstlisting}[language=Python, caption={Gemini client initialization}]
from google import genai

client = genai.Client(api_key="<GEMINI_API_KEY>")
response = client.models.generate_content(
    model="gemini-2.5-flash",
    contents=prompt
)
\end{lstlisting}

\textbf{Graceful Degradation:} All LLM endpoints are wrapped in try/except blocks. If the Gemini API is unreachable or rate-limited, the system returns a descriptive error message rather than crashing, ensuring the core deterministic pipeline continues to function independently.

%----------------------------------------------------------------------
\section{Feature 1: Empathetic Message Drafting}
\label{sec:draft-message}

\begin{table}[H]
\centering
\begin{tabular}{@{} l l @{}}
\toprule
\textbf{Property} & \textbf{Value} \\
\midrule
Endpoint     & \texttt{POST /api/llm/draft\_message} \\
Model        & \texttt{gemini-2.5-flash} \\
Input        & \texttt{MessageDraftRequest} (task, from, to, reason, benefits) \\
Output       & \texttt{\{"draft": "..."\}} \\
\bottomrule
\end{tabular}
\end{table}

When a manager accepts a workload reallocation suggestion from the AI Task Optimizer, they may need to communicate this change to the affected employee. This endpoint uses Gemini to draft a \textbf{polite, supportive, and non-confrontational} Slack/Email message.

\subsection{Prompt Design}

The prompt instructs the model to act as an \emph{empathetic management assistant} and incorporates:
\begin{itemize}
    \item The sender and recipient names
    \item The specific intervention (e.g., ``Workload Reallocation'', ``Schedule Focus Blocks'')
    \item The reason for the intervention (derived from burnout metrics)
    \item Expected benefits
    \item Strict tone constraints: supportive, blame-free, focused on well-being
\end{itemize}

\begin{lstlisting}[language=Python, caption={Message drafting prompt structure}]
prompt = (
    f"You are an empathetic and professional "
    f"management assistant.\n"
    f"Draft a polite, supportive, and "
    f"non-confrontational direct message "
    f"(Slack style) from '{from_name}' to "
    f"'{to_name}'.\n"
    f"The goal: suggest '{task}'.\n"
    f"Reason: {reason}\n"
    f"Benefits: {benefits}\n"
    f"Return only the drafted message text."
)
\end{lstlisting}

%----------------------------------------------------------------------
\section{Feature 2: Dynamic Employee Narratives}
\label{sec:llm-narrative}

\begin{table}[H]
\centering
\begin{tabular}{@{} l l @{}}
\toprule
\textbf{Property} & \textbf{Value} \\
\midrule
Endpoint     & \texttt{GET /api/llm/narrative/\{employee\_id\}} \\
Model        & \texttt{gemini-2.5-flash} \\
Input        & Employee ID (path parameter) \\
Output       & \texttt{\{"employeeId": "...", "narrative": "..."\}} \\
\bottomrule
\end{tabular}
\end{table}

This endpoint dynamically generates an insightful, empathetic narrative summarizing an employee's current well-being state. Unlike the rule-based narratives from \texttt{predictive\_ensemble.py} (Section~\ref{sec:ensemble}), these are generated in real-time by the LLM using the latest live metrics.

\subsection{Context Injection}

The following live metrics are injected into the prompt:
\begin{itemize}
    \item Burnout Probability (\%)
    \item Health Status (critical / warning / healthy)
    \item Deep Work Index (0--100)
    \item Fragmentation Score (0--100)
    \item Recovery Debt (hours)
\end{itemize}

The model is instructed to act as an \emph{organizational psychologist}, synthesize the numbers into insight (not just list them), identify one core challenge, and wrap with a supportive tone. The output is rendered in the Employee Detail view as an ``AI Risk Narrative'' card.

%----------------------------------------------------------------------
\section{Feature 3: Ask Pulse Copilot}
\label{sec:ask-pulse}

\begin{table}[H]
\centering
\begin{tabular}{@{} l l @{}}
\toprule
\textbf{Property} & \textbf{Value} \\
\midrule
Endpoint     & \texttt{POST /api/llm/ask\_pulse} \\
Model        & \texttt{gemini-2.5-flash} \\
Input        & \texttt{AskPulseRequest} (query, optional manager\_id) \\
Output       & \texttt{\{"response": "..."\}} \\
\bottomrule
\end{tabular}
\end{table}

The \textbf{Ask Pulse Copilot} is the most sophisticated Gen AI feature. It provides a \textbf{natural language conversational interface} within the Manager Dashboard, allowing managers to ask free-form questions about their team's burnout risks, workload distribution, and anomalies.

\subsection{RAG-Style Context Retrieval}

Ask Pulse simulates a \textbf{Retrieval-Augmented Generation (RAG)} architecture over PulseIQ's internal API:

\begin{enumerate}
    \item \textbf{Contextual data retrieval:} Before prompting the LLM, the endpoint calls internal functions to gather:
    \begin{itemize}
        \item All employees and their current metrics (filtered by \texttt{manager\_id} if provided)
        \item All pending workload intervention suggestions
    \end{itemize}

    \item \textbf{Context injection:} The gathered data is serialized and injected into the prompt as an ``\texttt{INTERNAL DATA CONTEXT}'' block.

    \item \textbf{Query resolution:} The manager's natural language question is appended, and Gemini answers using the provided real-time data.
\end{enumerate}

\begin{lstlisting}[language=Python, caption={Ask Pulse RAG-style prompt}]
# Gather real-time context
emps = get_employees(req.manager_id)
suggestions = get_suggestions(req.manager_id)
context_str = f"Employees: {emps}\n"
              f"Interventions: {suggestions}"

prompt = (
    f"You are 'Ask Pulse', the AI Copilot "
    f"for PulseIQ.\n"
    f"--- INTERNAL DATA CONTEXT ---\n"
    f"{context_str}\n"
    f"--- MANAGER'S QUERY ---\n"
    f"{req.query}\n"
)
\end{lstlisting}

\subsection{Frontend Integration}

The Ask Pulse Copilot is accessible as a dedicated navigation tab in the Manager Dashboard. It features:

\begin{itemize}
    \item A \textbf{chat-style interface} with user and AI message bubbles
    \item Real-time loading indicators (animated typing dots)
    \item Markdown rendering support for structured LLM responses
    \item Manager-scoped context (only shows data for the logged-in manager's reports)
\end{itemize}

Example queries a manager can ask:
\begin{quote}
    \emph{``Who on my team is most at risk of burnout right now?''}\\
    \emph{``What's the best way to redistribute Sarah's workload?''}\\
    \emph{``Are there any anomalies I should be aware of this week?''}
\end{quote}

%----------------------------------------------------------------------
\section{Pydantic Request Schemas}
\label{sec:genai-schemas}

The Gen AI endpoints use typed Pydantic models for input validation:

\begin{lstlisting}[language=Python, caption={Gen AI request schemas}]
class MessageDraftRequest(BaseModel):
    task: str          # e.g., "Workload Reallocation"
    from_name: str     # Manager name
    to_name: str       # Employee name
    reason: str        # Why the intervention
    benefits: List[str]# Expected outcomes

class AskPulseRequest(BaseModel):
    query: str                     # Free-form question
    manager_id: Optional[str] = None  # Scope context
\end{lstlisting}


%======================================================================
%  APPENDIX — Quick Reference
%======================================================================
\appendix
\chapter{Quick Reference}
\label{app:quickref}

\section{Project Directory Structure}

\begin{lstlisting}[language={}, caption={Key project directories and files}]
pulsefinal/
|-- run.py                    # Unified launcher
|-- backend/
|   |-- main.py               # Pipeline orchestrator
|   |-- api.py                # FastAPI server
|   |-- generate_data.py      # Synthetic data generator
|   |-- aggregate_features.py # Feature engineering
|   |-- compute_baselines.py  # Rolling baselines
|   |-- compute_burnout.py    # Behavioral engine
|   |-- recommendations.py    # Rule-based suggestions
|   |-- manager_insights.py   # Team rollups
|   |-- run_ml_parallel.py    # ML job coordinator
|   |-- time_series_analysis.py
|   |-- anomaly_detection.py
|   |-- deep_learning_model.py
|   |-- predictive_ensemble.py
|   |-- generate_report.py
|   `-- requirements.txt
|-- src/                      # React/Vite frontend
|-- pulseiq_data/             # Generated data artifacts
`-- Docs/                     # This documentation
\end{lstlisting}

\section{Generated Data Artifacts}

\begin{table}[H]
\centering
\caption{Key CSV artifacts produced by the pipeline}
\begin{tabular}{@{} l p{8cm} @{}}
\toprule
\textbf{File} & \textbf{Description} \\
\midrule
\texttt{*\_events.csv}            & Raw synthetic event data per source \\
\texttt{daily\_features.csv}      & Engineered features per (employee, date) \\
\texttt{daily\_scores.csv}        & Behavioral burnout scores \\
\texttt{timeseries\_forecasts.csv}& EWMA and trend forecasts \\
\texttt{anomaly\_scores.csv}      & Isolation Forest anomaly tags \\
\texttt{lstm\_predictions.csv}    & LSTM neural network predictions \\
\texttt{ensemble\_results.csv}    & Final combined predictions \\
\bottomrule
\end{tabular}
\end{table}

\section{Key Mathematical Formulas}

\begin{table}[H]
\centering
\caption{Summary of core mathematical formulas}
\begin{tabular}{@{} l l @{}}
\toprule
\textbf{Concept} & \textbf{Formula} \\
\midrule
Robust Z-Score        & $z = \dfrac{x - \tilde{x}}{\hat{\sigma}_{\text{robust}}}$ \\[10pt]
Fragmentation         & $F = 100 \cdot \tanh(\alpha \cdot z_{\text{agg}})$ \\[10pt]
Burnout Probability   & $P = \sigma(s) = \dfrac{1}{1 + e^{-s}}$ \\[10pt]
3-Day SMA             & $\text{SMA}_3(t) = \dfrac{P(t) + P(t\!-\!1) + P(t\!-\!2)}{3}$ \\[10pt]
EWMA                  & $E_t = \alpha x_t + (1-\alpha) E_{t-1}$ \\[10pt]
Zoom Speaking Ratio   & $r = \dfrac{\sum \text{speaking\_sec}}{\sum \text{duration}}$ \\
\bottomrule
\end{tabular}
\end{table}

%----------------------------------------------------------------------
\end{document}
